\documentclass[11pt,a4paper]{article}
\usepackage[utf8]{inputenc}
\usepackage[T1]{fontenc}
\usepackage[french]{babel}
\usepackage{amsmath,amsfonts,amssymb,bm}
\usepackage{booktabs}
\usepackage{microtype}
\usepackage[margin=2.5cm]{geometry}
\usepackage{hyperref}
\usepackage{listings}
\usepackage{xcolor}

\newcommand{\dd}{\mathrm{d}}

\definecolor{codegreen}{rgb}{0,0.6,0}
\definecolor{codegray}{rgb}{0.5,0.5,0.5}
\definecolor{codeblue}{rgb}{0.1,0.2,0.8}
\definecolor{backcolour}{rgb}{0.96,0.96,0.94}

\lstdefinestyle{mystyle}{
    backgroundcolor=\color{backcolour},
    commentstyle=\color{codegreen},
    keywordstyle=\color{codeblue},
    numberstyle=\tiny\color{codegray},
    stringstyle=\color{purple},
    basicstyle=\ttfamily\footnotesize,
    breaklines=true,
    keepspaces=true,
    showstringspaces=false,
    language=Python
}
\lstset{style=mystyle}

\title{Convergence du profil de curling $\omega(\rho)$\\[0.3em]
\large État des lieux, correctif du solveur BVP (commit \texttt{232334a})\\
et suppression de la minimisation de secours}
\author{JC Toussaint}
\date{12 octobre 2026}

\begin{document}
\maketitle

\section{Problème à résoudre}

Le profil d'aimantation du mode curling est décrit, en variable réduite
$\rho = r/R$, par l'équation d'Euler--Lagrange
\begin{equation}
  \omega'' = -\frac{\omega'}{\rho}
  + \left(\frac{1}{2\rho^{2}} + \kappa\right)\sin 2\omega,
  \qquad
  \kappa = \frac{K_u R^{2}}{2A},
  \label{eq:EL}
\end{equation}
avec les conditions aux limites
\begin{equation}
  \omega(\epsilon) = 0, \qquad \omega'(1) = 0, \qquad \epsilon = 10^{-6}.
\end{equation}
L'énergie associée est
\begin{equation}
  E[\omega] = 2\pi A \int_{\epsilon}^{1}
  \left[\rho\,\omega'^{2}
  + \left(\frac{1}{\rho} + 2\kappa\rho\right)\sin^{2}\omega\right]\dd\rho .
\end{equation}

Pour $K_u<0$ et $R>R_c$, $\omega$ croît de $0$ à $\pi/2$ dans une
\emph{couche limite} près du cœur, de largeur réduite
\begin{equation}
  \delta \sim \frac{\ell_K}{R} = \frac{1}{\sqrt{2|\kappa|}},
  \qquad \ell_K = \sqrt{A/|K_u|},
\end{equation}
puis reste à $\pi/2$ avec une pente nulle.

Cas test qui posait problème :
$A = 1.5\times10^{-11}$~J/m, $K_u = -10^{5}$~J/m$^{3}$, $R = 400$~nm,
d'où $\kappa \approx -533$ et $\delta \approx 0.03$.

\begin{lstlisting}[language=bash]
./curling_holo_simulator 1000 1e-9 --A 1.5e-11 --Ku=-1e5 --R 400e-9
\end{lstlisting}

\section{Stratégie précédente}

\begin{enumerate}
  \item Maillage uniforme : \texttt{rho\_grid = np.linspace(epsilon, 1, 600)}.
  \item Profil initial indépendant de $\kappa$ :
        $\omega_0(\rho) = \frac{\pi}{2}\tanh\!\big(\rho(2-\rho)\big)$.
  \item Résolution par \texttt{solve\_bvp} (\texttt{tol=1e-7}, \texttt{max\_nodes=1e6}).
  \item Test de monotonie sur $\omega'$. En cas d'échec, repli sur une
        minimisation directe de $E[\omega]$ :
        \begin{itemize}
          \item $\omega_i = \sum_{k\le i}\mathrm{softplus}(x_k)$, avec
                $\mathrm{softplus}(x) = \ln(1+e^{x}) > 0$, ce qui rend $\omega$
                croissant par construction ;
          \item optimisation par L-BFGS-B ;
          \item la condition $\omega'(1)=0$ n'est pas imposée : on compte sur
                le fait qu'elle est la condition naturelle du problème
                variationnel.
        \end{itemize}
\end{enumerate}

\subsection{Diagnostic sur le cas $\kappa=-533$}

\begin{center}
\begin{tabular}{@{}ll@{}}
\toprule
Étape & Résultat observé \\
\midrule
\texttt{solve\_bvp} & \emph{maximum number of mesh nodes exceeded}
                      (352\,615 nœuds) \\
                    & résidu $\sim 3\times10^{4}$,
                      $\omega$ oscille entre $-23$ et $+37$~rad \\
Repli (minimisation) & « convergé » en 40 itérations \\
                     & $\omega(1)-\pi/2 \approx +3.9\times10^{-3}$,
                       $\omega'(0.95)\approx 4\times10^{-2}$ \\
Temps total & 19.5~s \\
\bottomrule
\end{tabular}
\end{center}

\paragraph{Pourquoi le BVP divergeait.}
Le profil initial varie sur une échelle $\mathcal{O}(1)$ alors que la
solution varie sur $\delta\approx0.03$. Le maillage uniforme ne place qu'une
vingtaine de nœuds dans la couche limite. Newton part donc très loin de la
solution. Comme $\sin 2\omega$ est périodique, \eqref{eq:EL} admet d'autres
solutions (oscillantes, $\omega \to \omega + k\pi$), et le solveur s'égare
vers l'une d'elles en raffinant indéfiniment le maillage.

\paragraph{Pourquoi le repli ne donnait pas une pente nulle.}
C'est un défaut de structure de la paramétrisation par \emph{softplus} :
\begin{itemize}
  \item chaque incrément est strictement positif, donc le profil ne peut
        jamais être exactement plat ;
  \item une fois $\omega$ au-dessus de $\pi/2$, la contrainte de monotonie
        l'empêche de redescendre ;
  \item un incrément de l'ordre de $10^{-5}$ exige $x_k\approx-12$, où
        $\mathrm{softplus}'(x_k)\approx e^{x_k}\approx 10^{-5}$ : le gradient
        s'annule presque et L-BFGS-B ne peut plus réduire ces incréments.
\end{itemize}
Mettre la fonctionnelle à l'échelle (diviser par $2\pi A$) ne change rien :
le dépassement reste de $3\times10^{-3}$.
Enfin, la valeur $\omega'(1)\approx7\times10^{-6}$ renvoyée par la
\texttt{CubicSpline} est trompeuse : elle découle de la condition de bord
\emph{not-a-knot}, pas d'une vraie pente nulle (la pente de la dernière
maille vaut $\approx 1.1\times10^{-3}$).

\section{Correctif mis en œuvre}

Plutôt que de corriger le repli, on donne au solveur BVP les moyens de
converger, en tenant compte de l'échelle physique $\delta$.

\begin{lstlisting}
# Grid refined near the core (quadratic spacing) to resolve the boundary
# layer of width ~ l_K/R = 1/sqrt(2|kappa|), which is very thin for large |kappa|
rho_grid = epsilon + (1.0 - epsilon) * np.linspace(0.0, 1.0, N_BVP)**2
# 'guess' profile to push the solver toward the non-trivial solution
if Ku < 0:
    # The guess must vary on the physical length scale l_K/R, otherwise
    # solve_bvp diverges toward oscillating branches when |kappa| >> 1
    k = np.sqrt(2 * abs(kappa))
    omega_guess = (np.pi / 2) * np.tanh(k * rho_grid)
else:
    omega_guess = 0 * rho_grid
\end{lstlisting}

\paragraph{Profil initial à la bonne échelle.}
$\omega_0(\rho) = \frac{\pi}{2}\tanh\!\big(\sqrt{2|\kappa|}\,\rho\big)$
a déjà la forme de la solution : il monte sur une distance $\delta$, puis
sature à $\pi/2$. Il vérifie presque les conditions aux limites
($\omega_0(\epsilon)\approx0$ et $\omega_0'(1)\approx0$ dès que
$|\kappa|\gtrsim 1$). Newton démarre ainsi dans le bassin d'attraction de la
solution physique.

\paragraph{Maillage resserré près du cœur.}
Avec $\rho_i = \epsilon + (1-\epsilon)\,s_i^{2}$, $s_i$ uniforme sur $[0,1]$,
le pas près de $\epsilon$ vaut $\approx 3\times10^{-6}$, contre
$1.7\times10^{-3}$ en uniforme. Les nœuds se concentrent dans la couche
limite et autour de la singularité en $1/\rho^{2}$.

\paragraph{Cas limites.}
Pour $|\kappa|$ petit, $k$ est petit et le profil initial reste proche de
$0$, ce qui est cohérent avec $R<R_c$. Pour $K_u>0$, le profil initial reste
nul. L'interpolation (\texttt{sol.sol}) et le calcul de la phase ne
dépendent pas du maillage et n'ont pas été modifiés.

\section{Résultats}

\begin{center}
\begin{tabular}{@{}lrll@{}}
\toprule
Paramètres & $\kappa$ & Avant & Après \\
\midrule
$K_u=-10^{5}$, $R=400$~nm, $A=1.5\times10^{-11}$ & $-533$
  & divergence + repli imparfait & converge (2\,125 nœuds) \\
  & & 19.5~s & 0.47~s \\
Défaut ($K_u=-10^{4}$, $R=100$~nm) & $-3.85$
  & phase $\pm 8.8831$~rad & identique \\
$K_u=-2\times10^{4}$, $R=100$~nm, $A=1.5\times10^{-11}$ & $-6.67$
  & phase $\pm 4.9533$~rad & identique \\
$K_u=-10^{4}$, $R=20$~nm ($<R_c$) & $-0.15$
  & repli : $\omega(1)=0.586$ & $\omega\equiv0$ \\
  & & $E=+2.97\times10^{-11}$~J & $E\approx0$ \\
$K_u=+10^{4}$ & $+3.85$ & --- & converge \\
\bottomrule
\end{tabular}
\end{center}

Pour le cas $\kappa=-533$, la nouvelle solution vérifie
\[
  \omega'(1)\approx 10^{-31}, \qquad
  \omega(1)-\frac{\pi}{2}\approx -6\times10^{-15}, \qquad
  \omega' \ge 0 \text{ sur tout le domaine},
\]
et son énergie est plus basse que celle du repli :
\[
  \frac{E}{2\pi A} = -528.28144 \quad\text{(BVP)}
  \qquad\text{contre}\qquad
  -528.27107 \quad\text{(repli)} .
\]
Il s'agit donc d'un meilleur minimum. La phase holographique passe de
$\pm1.056$~rad à $\pm0.980$~rad.

Le cas $R<R_c$ est également corrigé : l'ancien repli produisait un profil
d'énergie strictement positive, donc supérieure à celle de la solution
triviale $\omega\equiv0$, que le BVP trouve désormais directement. La phase
passe de $\pm0.841$ à $\pm0.956$~rad, ce qui correspond à une aimantation
entièrement axiale ($\cos\omega = 1$).

\section{Suppression de la minimisation de secours}

\paragraph{Constat.}
Après le correctif, la minimisation n'était plus jamais appelée. Un balayage
de 224 jeux de paramètres, reprenant exactement les étapes de
\texttt{simulate\_curling()}, l'a vérifié :
\begin{center}
\begin{tabular}{@{}ll@{}}
\toprule
Paramètre & Valeurs \\
\midrule
$A$ (J/m) & $\{1.0,\ 1.3,\ 1.5,\ 2.0\}\times10^{-11}$ \\
$K_u$ (J/m$^{3}$) & $\{-10^{3},\ -10^{4},\ -5\times10^{4},\ -10^{5},\
                     -3\times10^{5},\ -10^{6},\ +10^{4},\ +10^{5}\}$ \\
$R$ (nm) & $\{10,\ 20,\ 50,\ 100,\ 200,\ 400,\ 800\}$ \\
\bottomrule
\end{tabular}
\end{center}
soit $|\kappa|$ entre $2.5\times10^{-3}$ et $3.2\times10^{4}$. Dans les 224
cas, \texttt{solve\_bvp} converge et la solution est croissante. Comme le repli
gardait ses défauts (section~2), il a été supprimé plutôt que conservé comme
code mort.

\paragraph{Code retiré.}
La section « 2 bis » (\texttt{total\_energy}, \texttt{minimize\_monotonic\_energy},
\texttt{MonotonicSplineWrapper}), les imports \texttt{scipy.optimize.minimize}
et \texttt{CubicSpline}, ainsi que le bloc de repli de
\texttt{simulate\_curling()}. Le calcul de $I(y)$ utilise directement
\texttt{sol.sol}.

\paragraph{Nouveau comportement en cas d'échec.}
Deux situations arrêtent désormais le programme avec un message d'erreur
sur la sortie d'erreur, la valeur de $\kappa$, et le code de sortie~1 :
\begin{itemize}
  \item \texttt{solve\_bvp} ne converge pas (\texttt{sol.status != 0}). Ce
        cas n'était pas testé auparavant : seule la monotonie l'était ;
  \item le profil n'est pas croissant ($\min\omega' < -10^{-12}$) ; le
        message indique alors la valeur minimale de $\omega'$.
\end{itemize}

\begin{lstlisting}
if sol.status != 0:
    print(f"{RED}[Euler-Lagrange / BVP] ERROR:{RESET} solve_bvp did not converge "
          f"(status {sol.status}: {sol.message})", file=sys.stderr)
    print(f"kappa = {kappa:.4g}: no physically admissible curling profile was found.",
          file=sys.stderr)
    sys.exit(1)

if np.all(domega_drho_bvp >= -1e-12):  # numerical tolerance
    print(f"{GREEN}[Monotonicity]{RESET} omega(rho) is increasing over the entire domain.")
else:
    print(f"{RED}[Monotonicity] ERROR:{RESET} omega(rho) is not monotonically increasing "
          f"over the domain (min domega/drho = {domega_drho_bvp.min():.3e}).", file=sys.stderr)
    print(f"kappa = {kappa:.4g}: the BVP solution does not satisfy the expected physical behavior.",
          file=sys.stderr)
    sys.exit(1)
\end{lstlisting}

\paragraph{Vérification.}
Les cas normaux donnent les mêmes phases qu'avec le repli présent et se
terminent avec le code~0. Un échec forcé (\texttt{max\_nodes=700} sur le cas
$\kappa=-533$) produit :
\begin{lstlisting}[language={}]
[Euler-Lagrange / BVP] ERROR: solve_bvp did not converge (status 1: The maximum number of mesh nodes is exceeded.)
kappa = -533.3: no physically admissible curling profile was found.
\end{lstlisting}
avec le code de sortie~1. Aucun jeu de paramètres réel testé n'atteint
l'erreur de monotonie : cette branche a seulement été vérifiée à la relecture.

\section{Points ouverts}

\begin{itemize}
  \item En cas d'échec, il n'y a plus de solution de repli. Si un besoin
        apparaît, une piste plus sûre que l'ancienne serait d'optimiser
        directement les valeurs $\omega_i$ avec des bornes
        $0\le\omega_i\le\pi/2$, sans paramétrisation \emph{softplus}.
  \item Pour des $|\kappa|$ encore plus grands, on pourrait ajouter une
        \emph{continuation en $\kappa$} : résoudre pour un $|\kappa|$ modéré,
        puis réutiliser la solution comme profil initial en augmentant
        $|\kappa|$ progressivement.
\end{itemize}

\end{document}
